% ============================================================================
%  Esquemas de compromiso y sus aplicaciones
%  Módulo 02 · Criptografía aplicada — ChainNotes
%
%  Versión LaTeX del apunte web temas/compromisos.html
%  Todas las figuras se generan con TikZ / pgfplots: no hay imágenes externas.
%
%  Compilar:  pdflatex esquemas-de-compromiso.tex   (dos pasadas, por el índice)
%             o bien:  latexmk -pdf esquemas-de-compromiso.tex
% ============================================================================

\documentclass[11pt,a4paper]{article}

% --- idioma y codificación -------------------------------------------------
\usepackage[T1]{fontenc}
\usepackage[utf8]{inputenc}
\usepackage[spanish,es-noshorthands,es-tabla]{babel}
\usepackage{lmodern}
\usepackage{microtype}

% --- página ----------------------------------------------------------------
\usepackage[a4paper,margin=2.6cm,headheight=14pt]{geometry}
\usepackage{fancyhdr}
\usepackage{enumitem}
\usepackage{booktabs}
\usepackage{tabularx}
\usepackage{array}
\usepackage{caption}

% --- matemáticas -----------------------------------------------------------
\usepackage{amsmath,amssymb,amsthm}
\usepackage{mathtools}

% --- gráficas --------------------------------------------------------------
\usepackage{tikz}
\usetikzlibrary{positioning,arrows.meta,calc,fit,backgrounds,shapes.geometric,
                decorations.pathreplacing}
\usepackage{pgfplots}
\pgfplotsset{compat=1.18}

% --- enlaces y cajas (hyperref antes que tcolorbox) ------------------------
\usepackage[colorlinks=true,linkcolor=violetdeep,urlcolor=violetdeep,
            citecolor=violetdeep,pdftitle={Esquemas de compromiso y sus aplicaciones},
            pdfauthor={ChainNotes}]{hyperref}
\usepackage[most]{tcolorbox}

% ---------------------------------------------------------------------------
%  Paleta (heredada de assets/css/styles.css)
% ---------------------------------------------------------------------------
\definecolor{violetc}{HTML}{7A68F5}
\definecolor{violetdeep}{HTML}{5B46E0}
\definecolor{mintc}{HTML}{1FA88A}
\definecolor{coralc}{HTML}{E04E5C}
\definecolor{amberc}{HTML}{C98A1F}
\definecolor{skyc}{HTML}{1E88C7}
\definecolor{inkc}{HTML}{282344}

\hypersetup{linkcolor=violetdeep}

% ---------------------------------------------------------------------------
%  Estilo de párrafo y encabezados
% ---------------------------------------------------------------------------
\setlength{\parindent}{0pt}
\setlength{\parskip}{0.55em}
\captionsetup{font=small,labelfont={bf,color=violetdeep}}

% Este documento lleva muchas figuras grandes. Con los valores por omisión LaTeX
% las aparta a páginas propias y deja huecos; aflojamos los umbrales.
\renewcommand{\topfraction}{0.9}
\renewcommand{\bottomfraction}{0.8}
\renewcommand{\textfraction}{0.07}
\renewcommand{\floatpagefraction}{0.75}
\setcounter{topnumber}{3}
\setcounter{totalnumber}{4}

\pagestyle{fancy}
\fancyhf{}
\fancyhead[L]{\small\color{inkc}\textsc{Esquemas de compromiso}}
\fancyhead[R]{\small\color{inkc}\thepage}
\renewcommand{\headrulewidth}{0.4pt}

% ---------------------------------------------------------------------------
%  Entornos de teoremas
% ---------------------------------------------------------------------------
\theoremstyle{definition}
\newtheorem{definicion}{Definición}
\theoremstyle{plain}
\newtheorem{teorema}{Teorema}
\newtheorem{proposicion}{Proposición}
\newtheorem{lema}{Lema}
\theoremstyle{remark}
\newtheorem*{observacion}{Observación}

% ---------------------------------------------------------------------------
%  Cajas de color
% ---------------------------------------------------------------------------
\newtcolorbox{nota}[1]{enhanced,breakable,colback=skyc!6,colframe=skyc!75!black,
  coltitle=white,fonttitle=\bfseries\small,title={#1},
  left=7pt,right=7pt,top=5pt,bottom=5pt,boxrule=0.7pt,arc=2.5pt}

\newtcolorbox{caso}[1]{enhanced,breakable,colback=amberc!8,colframe=amberc!85!black,
  coltitle=white,fonttitle=\bfseries\small,title={#1},
  left=7pt,right=7pt,top=5pt,bottom=5pt,boxrule=0.7pt,arc=2.5pt}

\newtcolorbox{aviso}[1]{enhanced,breakable,colback=coralc!7,colframe=coralc!85!black,
  coltitle=white,fonttitle=\bfseries\small,title={#1},
  left=7pt,right=7pt,top=5pt,bottom=5pt,boxrule=0.7pt,arc=2.5pt}

\newtcolorbox{laboratorio}[1]{enhanced,breakable,colback=violetc!6,colframe=violetc!80!black,
  coltitle=white,fonttitle=\bfseries\small,title={#1},
  left=7pt,right=7pt,top=5pt,bottom=5pt,boxrule=0.7pt,arc=2.5pt}

% ---------------------------------------------------------------------------
%  Atajos de notación
% ---------------------------------------------------------------------------
\newcommand{\bits}{\{0,1\}}
\newcommand{\concat}{\,\|\,}
\newcommand{\Adv}{\mathrm{Adv}}
\newcommand{\Com}{\mathsf{Com}}
\newcommand{\Open}{\mathsf{Open}}
\newcommand{\Setup}{\mathsf{Setup}}
\newcommand{\Zq}{\mathbb{Z}_q}

% ===========================================================================
\begin{document}
% ===========================================================================

% babel-spanish redefine \% inspeccionando \lastskip, lo que revienta dentro de
% los nodos de pgfplots. Restauramos la definición estándar de LaTeX.
\makeatletter
\DeclareRobustCommand{\%}{\char`\%\relax}
\makeatother

\begin{titlepage}
\centering
\vspace*{2.5cm}

{\Large\color{violetdeep}\textsc{ChainNotes}\par}
{\small\color{inkc}Módulo 02 · Criptografía aplicada\par}

\vspace{1.4cm}
{\Huge\bfseries Esquemas de compromiso\par}
\vspace{0.4cm}
{\LARGE y sus aplicaciones\par}

\vspace{1.2cm}
\begin{minipage}{0.78\textwidth}
\centering\itshape\color{inkc}
Un sobre cerrado que nadie puede leer y que tú ya no puedes cambiar. Con esa sola
idea se construyen las subastas honestas, los volados por teléfono, las pruebas
de conocimiento cero y las transacciones que se validan sin revelar cuánto dinero
llevan dentro.
\end{minipage}

\vspace{1.6cm}
\begin{tikzpicture}[x=6mm,y=6mm]
  \draw[violetdeep,line width=0.5pt,fill=violetc!12,rounded corners=0.6pt]
    (0,0) rectangle (5,3.2);
  \draw[violetdeep,line width=0.5pt] (0,3.2) -- (2.5,1.4) -- (5,3.2);
  \node[violetdeep,font=\small] at (2.5,0.9) {$c = \Com(m,r)$};
\end{tikzpicture}

\vfill
{\small Apuntes de Blockchain · Proyecto académico\par}
{\small\color{inkc}\today\par}
\end{titlepage}

\tableofcontents
\clearpage

% ===========================================================================
\section{Qué es comprometerse, exactamente}
% ===========================================================================

\begin{nota}{Objetivo del módulo}
Separar con precisión las dos propiedades de un compromiso ---ocultamiento y
vinculación--- y entender por qué cada una defiende contra un adversario
distinto.
\end{nota}

Dos personas que no se tienen confianza necesitan que una de ellas decida algo
\emph{ahora} y lo demuestre \emph{después}. Un compromiso es exactamente eso: un
sobre cerrado. Alice mete su decisión, lo entrega, y más tarde lo abre delante de
todos.

\begin{figure}[htbp]
\centering
\begin{tikzpicture}[
  font=\small,
  act/.style={draw=inkc!55,fill=inkc!8,rounded corners=2pt,
              minimum width=18mm,minimum height=6mm},
  msg/.style={draw=violetdeep,fill=violetc!12,rounded corners=2pt,inner sep=3pt,
              font=\scriptsize},
  >={Stealth[length=2mm]}
]
\node[act] (A) at (0,0) {Alice};
\node[act] (B) at (9,0) {Bob};
\draw[inkc!35,dashed] (A) -- ++(0,-5.0);
\draw[inkc!35,dashed] (B) -- ++(0,-5.0);

\node[msg,anchor=west] at (0.35,-1.0)
  {elige $m$, sortea $r \leftarrow \mathcal{R}$};
\draw[->,violetdeep,thick] (0,-2.0) -- node[above,font=\scriptsize,violetdeep]
  {$c = \Com(m, r)$} (9,-2.0);
\node[font=\scriptsize,skyc,anchor=north west,align=left] at (9.3,-2.1)
  {\textbf{ocultamiento:}\\no puede leer $m$};

\draw[->,mintc,thick] (0,-3.9) -- node[above,font=\scriptsize,mintc]
  {apertura: $(m, r)$} (9,-3.9);
\node[font=\scriptsize,coralc,anchor=north west,align=left] at (9.3,-4.0)
  {\textbf{vinculación:}\\$m$ ya no pudo cambiar};

\draw[decorate,decoration={brace,amplitude=4pt},inkc!55]
  (-1.5,-0.7) -- (-1.5,-2.3) node[midway,left=3pt,font=\scriptsize,inkc]{fase 1};
\draw[decorate,decoration={brace,amplitude=4pt},inkc!55]
  (-1.5,-3.6) -- (-1.5,-4.2) node[midway,left=3pt,font=\scriptsize,inkc]{fase 2};
\end{tikzpicture}
\caption{El protocolo, siempre en dos tiempos. Entre las dos fases puede pasar
cualquier cosa ---otras pujas, el desafío de un verificador, un bloque nuevo---
y el contenido del sobre sigue fijo y oculto.}
\label{fig:protocolo}
\end{figure}

El sobre tiene que cumplir dos cosas a la vez, y conviene verlas como dos
promesas a dos personas distintas:

\begin{description}[leftmargin=1.4em,style=unboxed,itemsep=0.3em]
  \item[Ocultamiento.] Es la promesa a Alice: quien tiene el sobre no puede ver
    lo que hay dentro. El adversario es \emph{el receptor}, que quiere leer antes
    de tiempo.
  \item[Vinculación.] Es la promesa a Bob: Alice ya no puede cambiar lo que
    metió. El adversario es \emph{el emisor}, que quiere cambiar de opinión
    cuando vea cómo van las cosas.
\end{description}

\begin{figure}[htbp]
\centering
\begin{tikzpicture}[
  font=\small,
  caja/.style={draw=#1,fill=#1!8,rounded corners=3pt,minimum width=58mm,
               minimum height=26mm,anchor=north west},
  >={Stealth[length=2.2mm]}
]
\node[caja=skyc] (p1) at (0,0) {};
\node[anchor=north west,font=\bfseries,skyc] at ($(p1.north west)+(3mm,-3mm)$)
  {Ataque al ocultamiento};
\node[anchor=north west,align=left,font=\scriptsize,inkc,text width=52mm]
  at ($(p1.north west)+(3mm,-9mm)$)
  {El \textbf{receptor} tiene $c$ y quiere deducir $m$ antes de la apertura.\\[2pt]
   Gana si distingue $\Com(m_0,r)$ de $\Com(m_1,r)$.};

\node[caja=coralc] (p2) at (6.6,0) {};
\node[anchor=north west,font=\bfseries,coralc] at ($(p2.north west)+(3mm,-3mm)$)
  {Ataque a la vinculación};
\node[anchor=north west,align=left,font=\scriptsize,inkc,text width=52mm]
  at ($(p2.north west)+(3mm,-9mm)$)
  {El \textbf{emisor} publicó $c$ y quiere abrirlo en otra cosa.\\[2pt]
   Gana si produce $(m,r) \neq (m',r')$ con el mismo $c$.};

\node[inkc,font=\scriptsize\itshape,anchor=north] at (6.3,-3.1)
  {dos adversarios opuestos: por eso un cifrado no basta};
\end{tikzpicture}
\caption{Las dos propiedades no son dos versiones de lo mismo. Protegen a partes
distintas, contra ataques distintos, y la sección~\ref{sec:imposibilidad}
demuestra que no pueden ser ambas perfectas.}
\label{fig:adversarios}
\end{figure}

\begin{aviso}{Compromiso no es cifrado}
Cifrar protege al que guarda el secreto. Comprometerse protege también
\emph{al otro}. Un esquema de cifrado no promete nada sobre la imposibilidad de
abrir el mismo criptograma de dos maneras; un compromiso vive o muere por eso.
\end{aviso}

\begin{definicion}[Esquema de compromiso]
Un esquema de compromiso es una terna de algoritmos $(\Setup, \Com, \Open)$ con
$\Com : \mathcal{M} \times \mathcal{R} \to \mathcal{C}$ tal que:
\begin{enumerate}[label=(\roman*),leftmargin=2.2em,itemsep=0.2em]
  \item \textbf{Corrección.} $\Open(c, m, r) = 1$ siempre que $c = \Com(m,r)$.
  \item \textbf{Ocultamiento.} Para $m_0 \neq m_1$ elegidos por el adversario,
    las distribuciones $\{\Com(m_0,r)\}_r$ y $\{\Com(m_1,r)\}_r$ con $r$ uniforme
    son indistinguibles:
    \[
      \Adv^{\text{ocult}}(\mathcal{A}) =
      \bigl|\Pr[\mathcal{A} \to 1 \mid m_0] - \Pr[\mathcal{A} \to 1 \mid m_1]\bigr|.
    \]
  \item \textbf{Vinculación.} Es inviable producir $(m,r) \neq (m',r')$ con
    $\Com(m,r) = \Com(m',r')$:
    \[
      \Adv^{\text{vinc}}(\mathcal{B}) =
      \Pr[\mathcal{B} \to \text{dos aperturas válidas de un mismo } c].
    \]
\end{enumerate}
Cada ventaja puede ser \emph{cero} (perfecta, incluso frente a un adversario sin
límite de cómputo) o \emph{despreciable} (computacional).
\end{definicion}

% ===========================================================================
\section{El compromiso de hash, y por qué $r$ no es opcional}
% ===========================================================================

El esquema más simple ya aparecía en el módulo 01: toma el mensaje, pégale un
valor aleatorio y hashea.
\begin{equation}\label{eq:hashcom}
  c = H(m \concat r), \qquad r \leftarrow \bits^{\lambda} \text{ uniforme}.
\end{equation}

De dónde sale cada propiedad es directo, y conviene decirlo con nombre y apellido
porque es el patrón que se repite en todo lo que sigue.

\begin{proposicion}\label{prop:hashbind}
La vinculación del esquema~\eqref{eq:hashcom} se reduce a la resistencia a
colisión de $H$: $\Adv^{\text{vinc}}(\mathcal{A}) \le \Adv^{\text{col}}_H(\mathcal{A})$.
\end{proposicion}

\begin{proof}
Sea $\mathcal{A}$ un adversario que produce $(m,r) \neq (m',r')$ con el mismo
compromiso. Entonces $H(m \concat r) = H(m' \concat r')$ con
$m \concat r \neq m' \concat r'$, es decir, una colisión de $H$ hallada
eficientemente.
\end{proof}

El ocultamiento, en cambio, \textbf{no viene de $H$}: viene de $r$. Éste es el
punto que casi todo el mundo se salta. La tentación es pensar que como $H$ es
irreversible, $c = H(m)$ ya oculta. Es falso, y de una manera brutal: el
adversario no necesita invertir nada. Le basta con hashear todos los mensajes
posibles y comparar. Si el espacio es «águila o sol», el ataque cuesta dos hashes.

Modelando $H$ como oráculo aleatorio, con $q$ consultas al oráculo,
\begin{equation}\label{eq:hidebound}
  \Adv^{\text{ocult}}(\mathcal{A}) \;\le\; \frac{q \cdot |\mathcal{M}|}{2^{\lambda}},
\end{equation}
y por eso $\lambda = 128$ basta para cualquier $\mathcal{M}$ realista. La
figura~\ref{fig:coste} pone la cota en escala.

\begin{figure}[htbp]
\centering
\begin{tikzpicture}
\begin{axis}[
  width=12.4cm, height=6.4cm,
  xlabel={bits de aleatoriedad en $r$ \quad ($\lambda$)},
  ylabel={$\log_2$ del trabajo del adversario},
  xmin=0, xmax=128, ymin=0, ymax=180,
  xtick={0,16,32,48,64,80,96,112,128},
  ytick={0,32,64,96,128,160},
  grid=major, grid style={inkc!12}, axis lines=left,
  tick label style={font=\footnotesize}, label style={font=\small},
  legend style={at={(0.03,0.97)},anchor=north west,font=\footnotesize,
                draw=inkc!30,fill=white},
  legend cell align=left
]
\fill[coralc!10] (axis cs:0,0) rectangle (axis cs:128,64);
\node[coralc,font=\scriptsize,anchor=west] at (axis cs:3,20)
  {zona alcanzable por un adversario bien financiado ($\le 2^{64}$)};
\addplot[skyc,very thick,domain=0:128,samples=2] {x+10};
\addlegendentry{$|\mathcal{M}| = 2^{10}$ (mil opciones)}
\addplot[amberc,very thick,domain=0:128,samples=2] {x+20};
\addlegendentry{$|\mathcal{M}| = 2^{20}$ (un millón)}
\addplot[mintc,very thick,domain=0:128,samples=2] {x+40};
\addlegendentry{$|\mathcal{M}| = 2^{40}$ (un billón)}
\draw[inkc!55,dashed] (axis cs:0,64) -- (axis cs:128,64);
\addplot[only marks,mark=*,mark size=2pt,violetdeep] coordinates {(0,10)(0,20)(0,40)};
\node[violetdeep,font=\scriptsize,anchor=west] at (axis cs:2,52)
  {$\lambda = 0$: el ataque cuesta sólo $|\mathcal{M}|$};
\end{axis}
\end{tikzpicture}
\caption{Coste del ataque al ocultamiento, $|\mathcal{M}|\cdot 2^{\lambda}$
hashes. El verificador honesto paga siempre \emph{un} hash, sea cual sea
$\lambda$: ésa asimetría es todo el esquema. Con $\lambda = 0$ el compromiso no
es más que un identificador público del mensaje.}
\label{fig:coste}
\end{figure}

\begin{laboratorio}{Laboratorio 2.1 — Cuánto cuesta abrir el sobre a la fuerza}
El adversario conoce el espacio de mensajes; lo que no conoce es $r$. La versión
interactiva ejecuta el ataque de verdad cuando cabe en un parpadeo, y cuando no,
mide la velocidad real de la máquina y extrapola.

\textbf{Qué observar.} Con $\lambda = 0$ y mil mensajes, el compromiso cae en
microsegundos. Con $\lambda = 128$, el mismo ataque pide $\sim 10^{41}$ hashes:
del orden de $10^{27}$ años en un portátil.
\end{laboratorio}

\begin{aviso}{El ocultamiento depende del espacio de mensajes}
Un hash es determinista: mensajes iguales dan compromisos iguales. Sólo cuando
$r$ es genuinamente aleatorio y suficientemente largo el conjunto de candidatos
se vuelve impracticable de recorrer. Ni el ocultamiento ni la vinculación de este
esquema son perfectos: un adversario sin límite de cómputo rompe ambos.
\end{aviso}

% ===========================================================================
\section{No puedes tener las dos cosas perfectas}
\label{sec:imposibilidad}
% ===========================================================================

Antes de buscar esquemas mejores conviene saber qué es imposible, porque ahorra
mucho tiempo.

\begin{proposicion}\label{prop:imposible}
Ningún esquema de compromiso es simultáneamente perfectamente ocultante y
perfectamente vinculante.
\end{proposicion}

\begin{proof}
Supongamos ocultamiento perfecto. Para cualesquiera $m_0 \neq m_1$, las dos
distribuciones
\[
  \{\Com(m_0,r)\}_{r} \quad\text{y}\quad \{\Com(m_1,r)\}_{r}
\]
son idénticas; en particular tienen el mismo soporte. Luego para todo $c$ en el
soporte y todo $m \in \mathcal{M}$ existe $r$ con $\Com(m,r) = c$.

Un emisor sin límite de cómputo enumera $\mathcal{R}$, encuentra ese $r$ y abre
$c$ donde quiera. La vinculación perfecta falla.

El recíproco es simétrico: si la vinculación es perfecta, cada $c$ admite un
único $m$, así que $c$ determina $m$ y un adversario sin límites lo recupera por
búsqueda exhaustiva.
\end{proof}

\begin{figure}[htbp]
\centering
\begin{tikzpicture}[
  font=\small,
  celda/.style={draw=#1,fill=#1!8,rounded corners=3pt,minimum width=56mm,
                minimum height=24mm,anchor=north west},
  muerta/.style={draw=inkc!35,fill=inkc!5,rounded corners=3pt,minimum width=56mm,
                 minimum height=24mm,anchor=north west}
]
\node[celda=skyc] (c1) at (0,0) {};
\node[anchor=north west,font=\bfseries,skyc] at ($(c1.north west)+(3mm,-3mm)$)
  {Ocultamiento perfecto};
\node[anchor=north west,align=left,font=\scriptsize,inkc,text width=50mm]
  at ($(c1.north west)+(3mm,-8.5mm)$)
  {El sobre no filtra nada, ni ante cómputo ilimitado. La vinculación se apoya en
   un problema difícil.};
\node[anchor=south west,font=\scriptsize\bfseries,skyc]
  at ($(c1.south west)+(3mm,2mm)$) {PEDERSEN};

\node[celda=coralc] (c2) at (6.4,0) {};
\node[anchor=north west,font=\bfseries,coralc] at ($(c2.north west)+(3mm,-3mm)$)
  {Vinculación perfecta};
\node[anchor=north west,align=left,font=\scriptsize,inkc,text width=50mm]
  at ($(c2.north west)+(3mm,-8.5mm)$)
  {Cada compromiso admite una sola apertura, siempre. El secreto sólo dura
   mientras el adversario sea finito.};
\node[anchor=south west,font=\scriptsize\bfseries,coralc]
  at ($(c2.south west)+(3mm,2mm)$) {ELGAMAL COMO COMPROMISO};

\node[celda=mintc] (c3) at (0,-2.9) {};
\node[anchor=north west,font=\bfseries,mintc] at ($(c3.north west)+(3mm,-3mm)$)
  {Ambas computacionales};
\node[anchor=north west,align=left,font=\scriptsize,inkc,text width=50mm]
  at ($(c3.north west)+(3mm,-8.5mm)$)
  {Lo que da el compromiso de hash. Suficiente en la práctica y el más barato de
   todos: un hash y nada más.};
\node[anchor=south west,font=\scriptsize\bfseries,mintc]
  at ($(c3.south west)+(3mm,2mm)$) {$H(m \concat r)$};

\node[muerta] (c4) at (6.4,-2.9) {};
\node[anchor=north west,font=\bfseries,inkc!60] at ($(c4.north west)+(3mm,-3mm)$)
  {Ambas perfectas};
\node[anchor=north west,align=left,font=\scriptsize,inkc!60,text width=50mm]
  at ($(c4.north west)+(3mm,-8.5mm)$)
  {Imposible. No hay ingeniería que lo arregle: es un argumento de conteo, no una
   limitación tecnológica.};
\node[anchor=south west,font=\scriptsize\bfseries,inkc!60]
  at ($(c4.south west)+(3mm,2mm)$) {NO EXISTE};
\draw[coralc,line width=1pt] ($(c4.north west)+(2mm,-2mm)$) -- ($(c4.south east)+(-2mm,2mm)$);
\end{tikzpicture}
\caption{El cuadrante completo. Tres casillas están habitadas y la cuarta está
vacía por la Proposición~\ref{prop:imposible}.}
\label{fig:cuadrante}
\end{figure}

\begin{nota}{La pregunta de diseño}
No es «¿cuál es más seguro?», sino \textbf{¿qué prefieres que sobreviva a un
adversario del futuro?} Si el secreto debe seguir oculto en cincuenta años, elige
ocultamiento perfecto. Si lo que no puede cambiar nunca es el registro, elige
vinculación perfecta.
\end{nota}

% ===========================================================================
\section{Pedersen: el compromiso que además suma}
% ===========================================================================

\begin{nota}{Objetivo del módulo}
Construir el compromiso de Pedersen en un grupo concreto, demostrar que su
ocultamiento es perfecto y que su vinculación equivale al logaritmo discreto.
\end{nota}

El compromiso de hash funciona, pero es opaco: una vez cerrado el sobre no puedes
hacer \emph{nada} con él salvo abrirlo. Pedersen cambia eso, y el precio que paga
es necesitar estructura algebraica en vez de una función de mezcla.

\begin{definicion}[Compromiso de Pedersen]
Sea $G$ un grupo cíclico de orden primo $q$ donde el logaritmo discreto es
difícil, y sean $g, h \in G$ dos generadores tales que nadie conoce
$\log_g h$. Para $m, r \in \Zq$ con $r$ uniforme:
\[
  C(m, r) = g^{m} \cdot h^{r}.
\]
\end{definicion}

En los laboratorios de este apunte $G$ es el subgrupo de residuos cuadráticos
módulo un primo seguro $p = 2q+1$. Es un grupo real, no una maqueta: la
aritmética es la misma que usarías en producción, sólo que con números pequeños
para que quepan en pantalla y para que los ataques se puedan ejecutar.

\begin{proposicion}[Ocultamiento perfecto]\label{prop:pedhide}
El compromiso de Pedersen es perfectamente ocultante: para todo $m$, la
distribución de $C(m,r)$ con $r$ uniforme en $\Zq$ es la uniforme sobre $G$.
\end{proposicion}

\begin{proof}
Como $h$ genera $G$ y $r$ es uniforme en $\Zq$, el elemento $h^{r}$ es uniforme
sobre $G$. Multiplicar por la constante $g^{m}$ es una biyección de $G$ en $G$,
y una biyección aplicada a una distribución uniforme da la uniforme. Luego
$C(m,r)$ es uniforme sobre $G$ \emph{independientemente de $m$}, y las
distribuciones para $m_0$ y $m_1$ son idénticas: $\Adv^{\text{ocult}} = 0$.
\end{proof}

\begin{figure}[htbp]
\centering
\begin{tikzpicture}
\begin{axis}[
  width=12.4cm, height=5.6cm,
  xlabel={tramos del grupo $G$ (el rango $[0,p)$ dividido en 48 partes)},
  ylabel={muestras},
  xmin=0, xmax=49, ymin=0, ymax=150,
  xtick={1,8,16,24,32,40,48},
  grid=major, grid style={inkc!12}, axis lines=left,
  tick label style={font=\footnotesize}, label style={font=\small},
  legend style={at={(0.5,-0.26)},anchor=north,legend columns=2,
                font=\footnotesize,draw=inkc!30,fill=white,
                /tikz/every even column/.append style={column sep=12pt}}
]
\addplot[ybar,bar width=4.4pt,fill=violetc!45,draw=violetc!80!black] coordinates {
  (1,97) (2,104) (3,98) (4,96) (5,112) (6,93) (7,98) (8,103) (9,85) (10,120)
  (11,80) (12,122) (13,98) (14,96) (15,105) (16,100) (17,89) (18,100) (19,86)
  (20,107) (21,86) (22,96) (23,83) (24,108) (25,106) (26,101) (27,102) (28,125)
  (29,97) (30,102) (31,105) (32,100) (33,92) (34,109) (35,96) (36,105) (37,79)
  (38,99) (39,98) (40,101) (41,99) (42,105) (43,122) (44,105) (45,107) (46,84)
  (47,100) (48,99)};
\addlegendentry{4\,800 compromisos de $m = 42$, con $r$ al azar}
\addplot[mintc,very thick,domain=0:49,samples=2] {100};
\addlegendentry{esperanza uniforme: $4800/48 = 100$}
\end{axis}
\end{tikzpicture}
\caption{Ocultamiento perfecto, dibujado. Son 4\,800 compromisos del
\emph{mismo} mensaje $m = 42$ en el grupo $q = 100\,043$, agrupados por tramo del
grupo. La dispersión observada (mínimo 79, máximo 125) es exactamente el ruido
multinomial esperado, $\sigma \approx 10$. El compromiso de un mensaje fijo puede
caer en cualquier parte con igual probabilidad.}
\label{fig:dispersion}
\end{figure}

\begin{proposicion}[Vinculación computacional]\label{prop:pedbind}
Si el logaritmo discreto es difícil en $G$, el compromiso de Pedersen es
computacionalmente vinculante. Más aún, toda ruptura de la vinculación
\emph{entrega} el logaritmo discreto.
\end{proposicion}

\begin{proof}
Supongamos que $\mathcal{A}$ produce dos aperturas válidas del mismo compromiso,
\[
  g^{m}h^{r} = g^{m'}h^{r'}, \qquad (m,r) \neq (m',r').
\]
Reagrupando, $g^{m-m'} = h^{r'-r}$. Nótese que $r \neq r'$: si fueran iguales
tendríamos $g^{m} = g^{m'}$ y por tanto $m = m'$, contradiciendo la hipótesis.
Como $q$ es primo y $0 < |r'-r| < q$, el valor $r'-r$ es invertible módulo $q$, y
\[
  \log_g h \;=\; (m - m')\,(r' - r)^{-1} \bmod q .
\]
Así que $\Adv^{\text{vinc}}(\mathcal{A}) \le \Adv^{\text{dlog}}(\mathcal{B})$.
\end{proof}

\begin{laboratorio}{Laboratorios 2.2 y 2.3 — Pedersen en vivo, y toda apertura es válida}
El primero calcula compromisos en el grupo elegido y dibuja la
figura~\ref{fig:dispersion} en tiempo real. El segundo toma un compromiso
publicado y, para cada mensaje de una lista arbitraria, busca el $r'$ que lo
abriría usando paso de bebé, paso de gigante.

\textbf{Qué observar.} Con $q = 100\,043$ los encuentra \emph{todos} en unos
2\,600 pasos y tres milisegundos. Las dos lecturas del mismo experimento: que
\emph{exista} un $r'$ para cada mensaje es el ocultamiento perfecto; que se haya
podido \emph{calcular} es la vinculación rota, y sólo ocurre porque $\sqrt{q}$
son unos cientos de operaciones. Con $q$ de 255 bits harían falta $2^{127}$
operaciones y $2^{127}$ palabras de memoria.
\end{laboratorio}

\begin{nota}{En producción}
Nadie usa enteros módulo $p$ para esto: se usan curvas elípticas ---secp256k1,
ristretto255, BLS12-381---, donde un elemento del grupo ocupa 32 o 48 bytes en
vez de 256. La aritmética cambia de nombre pero el esquema es literalmente el
mismo.
\end{nota}

% ===========================================================================
\section{Aritmética sobre sobres cerrados}
% ===========================================================================

Aquí aparece lo que el compromiso de hash no puede hacer. Multiplica dos
compromisos de Pedersen y mira qué sale:
\begin{equation}\label{eq:homo}
  C(a, r_1) \cdot C(b, r_2)
  = g^{a+b}\,h^{r_1+r_2}
  = C(a+b,\; r_1+r_2).
\end{equation}

Es una consecuencia de las leyes de los exponentes, casi decepcionante de lo
simple que es. Y sin embargo abre una puerta enorme: \textbf{se puede verificar
una suma sin conocer los sumandos}.

\begin{figure}[htbp]
\centering
\begin{tikzpicture}[
  font=\small,
  sobre/.style={draw=violetdeep,fill=violetc!10,rounded corners=2pt,
                minimum width=26mm,minimum height=11mm,align=center,font=\footnotesize},
  claro/.style={draw=inkc!45,fill=white,rounded corners=2pt,
                minimum width=20mm,minimum height=8mm,align=center,font=\footnotesize},
  >={Stealth[length=2.2mm]}
]
% mundo abierto
\node[claro] (a) at (0,1.6) {$a = 1250$};
\node[claro] (b) at (0,0) {$b = 3400$};
\node[claro] (s) at (0,-1.9) {$a+b = 4650$};
\draw[->,inkc!55] (a.south) -- (s.north west);
\draw[->,inkc!55] (b.south) -- (s.north);
\node[inkc,font=\scriptsize,anchor=east,align=right] at (-1.3,0.8) {suma\\ordinaria};

% flechas Com
\node[sobre] (ca) at (6.4,1.6) {$C(a, r_1)$};
\node[sobre] (cb) at (6.4,0) {$C(b, r_2)$};
\node[sobre] (cs) at (6.4,-1.9) {$C(a{+}b,\, r_1{+}r_2)$};
\draw[->,violetdeep] (a) -- node[above,font=\scriptsize,violetdeep]{$\Com$} (ca);
\draw[->,violetdeep] (b) -- node[above,font=\scriptsize,violetdeep]{$\Com$} (cb);
\draw[->,violetdeep] (s) -- node[above,font=\scriptsize,violetdeep]{$\Com$} (cs);

\draw[->,mintc,thick] (ca.south) -- (cs.north west);
\draw[->,mintc,thick] (cb.south) -- (cs.north);
\node[mintc,font=\scriptsize,anchor=west,align=left] at (7.9,0.8)
  {producto\\en el grupo};

\node[inkc,font=\scriptsize\itshape,anchor=north,align=center] at (3.2,-2.9)
  {el diagrama conmuta: da igual sumar y luego cerrar el sobre,\\
   que cerrar los sobres y luego multiplicarlos};
\end{tikzpicture}
\caption{El homomorfismo de la ecuación~\eqref{eq:homo}. El verificador recorre
el camino de la derecha, que sólo usa datos públicos, y concluye algo sobre el
camino de la izquierda, que nunca ve.}
\label{fig:homomorfismo}
\end{figure}

\subsection{Transacciones confidenciales}

De ahí sale la aplicación estrella. Una transacción es válida si lo que entra es
igual a lo que sale más la comisión. Si cada monto se publica como compromiso de
Pedersen y el emisor elige los factores de cegado de modo que se cancelen
---es decir, $\sum r_{\text{ent}} = \sum r_{\text{sal}} + r_{\text{com}}$---, el
nodo puede comprobar
\begin{equation}\label{eq:balance}
  \prod_i C_{\text{ent},i} \;=\; \prod_j C_{\text{sal},j} \cdot C_{\text{com}}
\end{equation}
que es una igualdad entre elementos del grupo, sin un solo monto a la vista.
Es exactamente lo que hacen Elements y Liquid con las \emph{confidential
transactions}, y lo que Mimblewimble ---Grin, Beam--- lleva al extremo
convirtiendo la cadena entera en una suma de compromisos.

\begin{figure}[htbp]
\centering
\begin{tikzpicture}[
  font=\footnotesize,
  ent/.style={draw=skyc,fill=skyc!10,rounded corners=2pt,
              minimum width=24mm,minimum height=9mm,align=center},
  sal/.style={draw=mintc,fill=mintc!10,rounded corners=2pt,
              minimum width=24mm,minimum height=9mm,align=center},
  com/.style={draw=amberc,fill=amberc!14,rounded corners=2pt,
              minimum width=20mm,minimum height=9mm,align=center},
  >={Stealth[length=2mm]}
]
\node[ent] (e1) at (0,1.1)  {$C_{\text{ent},1}$};
\node[ent] (e2) at (0,-0.5) {$C_{\text{ent},2}$};
\node[draw=violetdeep,fill=violetc!12,circle,inner sep=3pt] (prod) at (3.4,0.3) {$\prod$};
\node[draw=violetdeep,fill=violetc!12,circle,inner sep=3pt] (prod2) at (7.4,0.3) {$\prod$};
\node[sal] (s1) at (10.8,1.5)  {$C_{\text{sal},1}$};
\node[sal] (s2) at (10.8,0.0)  {$C_{\text{sal},2}$};
\node[com] (cm) at (10.8,-1.5) {$C_{\text{com}}$};

\draw[->,skyc] (e1) -- (prod); \draw[->,skyc] (e2) -- (prod);
\draw[->,mintc] (s1) -- (prod2); \draw[->,mintc] (s2) -- (prod2);
\draw[->,amberc] (cm) -- (prod2);
\draw[violetdeep,very thick] (prod) -- node[above,font=\small,violetdeep]{$\stackrel{?}{=}$} (prod2);

\node[inkc,font=\scriptsize,anchor=north,align=center] at (5.4,-1.2)
  {el nodo compara dos elementos del grupo\\y acepta o rechaza el bloque};
\end{tikzpicture}
\caption{Verificación de balance de la ecuación~\eqref{eq:balance}. Nada de lo
que el nodo toca es un monto.}
\label{fig:confidencial}
\end{figure}

\begin{aviso}{Y aquí está la trampa: el desbordamiento}
La igualdad~\eqref{eq:balance} se comprueba \textbf{módulo $q$}. Nada impide que
un atacante use montos gigantescos que den la vuelta al grupo. En el grupo del
laboratorio, con $q = 100\,043$, una transacción con entradas
$3 + 2 = 5$ y salidas $100\,038 + 10 = 100\,048 \equiv 5 \pmod q$ \emph{cuadra
perfectamente} y el nodo la acepta, habiendo creado cien mil unidades de la nada.

La defensa no es otro compromiso: son las \textbf{pruebas de rango}, que
demuestran en cero conocimiento que cada monto vive en $[0, 2^{64})$ y que la
suma no da la vuelta. Bulletproofs y sus sucesores existen exactamente para
tapar este agujero, y son la parte cara de una transacción confidencial: el
compromiso ocupa 32 bytes; la prueba de rango, cientos.
\end{aviso}

\begin{observacion}[Qué da y qué no da el homomorfismo]
$\Com$ es un homomorfismo de $(\Zq\times\Zq, +)$ en $(G, \cdot)$, lo que regala
la resta ($C(a)\,C(b)^{-1} = C(a-b)$) y el escalado por constante pública
($C(a)^{k} = C(ka)$). Lo que \emph{no} da es el producto: no existe forma de
obtener $C(ab)$ a partir de $C(a)$ y $C(b)$. Para eso hacen falta
emparejamientos o cifrado completamente homomórfico.
\end{observacion}

\begin{laboratorio}{Laboratorios 2.4 y 2.5 — Sumar sin abrir; transacción confidencial}
El primero compromete dos montos con factores de cegado al azar, multiplica los
compromisos y lo compara con el compromiso directo de la suma: coinciden siempre.
El segundo monta una transacción de dos entradas, dos salidas y comisión, con
tres escenarios preconfigurados.

\textbf{Qué observar.} El escenario «crear dinero» se rechaza, como debe. El
escenario «desbordamiento» se \textbf{acepta} y sin embargo crea dinero: es la
demostración ejecutable de por qué las pruebas de rango no son opcionales.
\end{laboratorio}

% ===========================================================================
\section{La puerta trasera que cabe en un solo número}
% ===========================================================================

Volvamos a la fase que parecía un trámite. El compromiso de Pedersen necesita dos
generadores, $g$ y $h$. La pregunta incómoda: ¿quién los eligió?

\begin{lema}[Equivocación con trampa]\label{lema:trapdoor}
Si $h = g^{x}$ con $x$ conocido, entonces para todo $c = C(m,r)$ y todo
$m' \in \Zq$ el valor
\[
  r' = r + (m - m')\,x^{-1} \bmod q
\]
satisface $C(m', r') = c$.
\end{lema}

\begin{proof}
Con $h = g^{x}$ todo compromiso colapsa a una sola potencia de $g$:
$C(m,r) = g^{m}g^{xr} = g^{m+xr}$. Así que $C(m',r') = C(m,r)$ equivale a
$m' + xr' \equiv m + xr \pmod q$, una ecuación lineal en $r'$ cuya solución es
la del enunciado; $x$ es invertible porque $q$ es primo y $x \neq 0$.
\end{proof}

Fíjate en lo que significa: la vinculación del esquema no depende sólo de que el
logaritmo discreto sea difícil, sino de que \textbf{nadie sepa ya la respuesta}
para ese par concreto. Ese $x$ se llama \emph{trampa} o \emph{trapdoor}.

\begin{figure}[htbp]
\centering
\begin{tikzpicture}[
  font=\footnotesize,
  ap/.style={draw=#1,fill=#1!10,rounded corners=2pt,
             minimum width=32mm,minimum height=10mm,align=center},
  >={Stealth[length=2.2mm]}
]
\node[draw=violetdeep,fill=violetc!14,rounded corners=3pt,
      minimum width=34mm,minimum height=12mm,align=center] (c) at (0,0)
  {$c$ publicado\\[1pt]\scriptsize un solo elemento del grupo};

\node[ap=mintc] (o1) at (6.6,1.4) {apertura $(m, r)$};
\node[ap=coralc] (o2) at (6.6,-1.4) {apertura $(m', r')$};

\draw[->,mintc,thick] (c) -- node[above,font=\scriptsize,mintc,sloped]
  {la honesta} (o1);
\draw[->,coralc,thick] (c) -- node[below,font=\scriptsize,coralc,sloped]
  {con $x$, en microsegundos} (o2);

\node[anchor=west,align=left,font=\scriptsize,inkc,text width=46mm] at (8.9,0)
  {Ambas verifican. El verificador no tiene forma de saber cuál era la
   «verdadera»: son indistinguibles por construcción.};
\end{tikzpicture}
\caption{Quien conoce la trampa rompe la \emph{vinculación}, no el ocultamiento.
Sigue sin poder leer sobres ajenos; lo que puede es mentir sobre los propios. En
una subasta eso basta para ganar siempre.}
\label{fig:trapdoor}
\end{figure}

\subsection{Las dos defensas}

\begin{description}[leftmargin=1.4em,style=unboxed,itemsep=0.35em]
  \item[Nothing up my sleeve.] Derivar $h$ de un hash de una cadena pública y sin
    gracia ---el nombre del protocolo, los dígitos de $\pi$---. Nadie conoce el
    logaritmo porque nadie eligió el resultado. Es lo que hacen los laboratorios
    de este apunte:
    \[
      h = \bigl(\mathrm{SHA\text{-}256}(\texttt{"ChainNotes/pedersen/h/0"})
          \bmod p\bigr)^{2} \bmod p .
    \]
  \item[Ceremonia multiparte.] Cuando los parámetros son demasiado complejos para
    derivarlos de un hash, muchos participantes aportan aleatoriedad por turnos.
    La trampa sólo existe si \emph{todos} conspiraron.
\end{description}

\begin{caso}{El «residuo tóxico»}
En el argot de las ceremonias, la aleatoriedad que cada participante debe
destruir se llama \emph{toxic waste}. Si sobrevive una sola copia, quien la tenga
puede falsificar. Zcash montó en 2016 una ceremonia con hardware desconectado y
actas públicas justamente por eso; la ceremonia de EIP-4844 reunió más de
140\,000 participantes con la misma lógica.
\end{caso}

\begin{laboratorio}{Laboratorio 2.6 — Abrir el mismo sobre de dos maneras}
Alice se compromete a $m$ y después decide que prefería $m'$. Con la trampa lo
consigue sin tocar el compromiso publicado, aplicando el
Lema~\ref{lema:trapdoor}.

\textbf{Qué observar.} Al activar «$h$ derivado de un hash», el laboratorio ya no
tiene $x$ y la fórmula deja de existir: cambiar de opinión pasa a exigir un
logaritmo discreto. El mismo esquema, el mismo código, y la seguridad depende
por completo de cómo se eligió un número.
\end{laboratorio}

% ===========================================================================
\section{Comprometerse a un millón de cosas a la vez}
% ===========================================================================

Hasta aquí el sobre guardaba un valor. ¿Y si hay que comprometerse a una lista
entera ---todas las transacciones de un bloque, el estado completo de una
cadena--- y además poder demostrar después que un elemento concreto estaba
dentro? Eso es un \emph{compromiso vectorial}, y del módulo 01 ya conoces uno: el
árbol de Merkle. La raíz es el compromiso; la prueba de inclusión es la apertura
\emph{parcial}.

\subsection{Compromisos polinómicos: KZG}

La idea de KZG es tratar la lista como los coeficientes de un polinomio $f$ y
comprometerse a $f$ evaluándolo «dentro del exponente» en un punto secreto $\tau$
que nadie conoce ---de ahí la ceremonia de la sección anterior---:
\[
  C = g^{f(\tau)}.
\]
Lo asombroso es la prueba. Para convencer a alguien de que $f(z) = y$ basta con
un elemento más del grupo: el compromiso al cociente
\[
  \pi = g^{\,\psi(\tau)}, \qquad
  \psi(X) = \frac{f(X) - y}{X - z},
\]
que sólo es un polinomio si efectivamente $f(z) = y$. No crece con el tamaño de
la lista: \textbf{48 bytes, siempre}. Ésa es la razón de que Ethereum haya
adoptado KZG para los \emph{blobs} de EIP-4844.

\begin{figure}[htbp]
\centering
\begin{tikzpicture}
\begin{axis}[
  width=12.4cm, height=6.4cm,
  ymode=log, log basis y=10,
  xlabel={elementos comprometidos ($n = 2^{k}$)},
  ylabel={bytes para probar un elemento},
  xmin=1, xmax=25.5, ymin=8, ymax=3e9,
  xtick={1,4,8,12,16,20,24},
  xticklabel={$2^{\pgfmathprintnumber{\tick}}$},
  grid=major, grid style={inkc!12}, axis lines=left,
  tick label style={font=\footnotesize}, label style={font=\small},
  legend style={at={(0.5,-0.26)},anchor=north,legend columns=1,
                font=\footnotesize,draw=inkc!30,fill=white},
  legend cell align=left
]
\addplot[coralc,very thick,domain=1:24,samples=24] {2^x*32};
\addlegendentry{mandar la lista completa: $n \cdot 32$ B}
\addplot[amberc,very thick,domain=1:24,samples=24] {x*32};
\addlegendentry{prueba de Merkle: $\lceil\log_2 n\rceil \cdot 32$ B}
\addplot[mintc,very thick,domain=1:24,samples=2] {48};
\addlegendentry{prueba KZG: 48 B, constante}
\addplot[only marks,mark=*,mark size=2pt,inkc] coordinates {(20,33554432)(20,640)(20,48)};
\node[inkc,font=\scriptsize,anchor=south east] at (axis cs:19.6,33554432) {32 MB};
\node[inkc,font=\scriptsize,anchor=south east] at (axis cs:19.6,640) {640 B};
\node[inkc,font=\scriptsize,anchor=south east] at (axis cs:19.6,48) {48 B};
\draw[violetdeep,dashed] (axis cs:20,8) -- (axis cs:20,3e9);
\node[violetdeep,font=\scriptsize,anchor=north east,align=right]
  at (axis cs:19.6,2.4e9) {un millón\\de elementos};
\end{axis}
\end{tikzpicture}
\caption{Con un millón de elementos la diferencia es de 32 MB a 640 bytes a 48
bytes. El ahorro de Merkle es logarítmico; el de KZG es total, y por eso se paga
una ceremonia por él.}
\label{fig:tamanos}
\end{figure}

\begin{table}[htbp]
\centering\small
\caption{Las tres familias de compromisos a muchos valores.}
\label{tab:vectoriales}
\begin{tabularx}{\textwidth}{@{}l l l >{\raggedright\arraybackslash}X@{}}
\toprule
\textbf{Familia} & \textbf{Compromiso} & \textbf{Apertura parcial} &
\textbf{Qué se paga} \\
\midrule
Merkle & 32 B & $\lceil\log_2 n\rceil \cdot 32$ B &
  Nada: sólo hashes, sin setup ni supuestos algebraicos. Resistente a cuántico. \\
\addlinespace
Pedersen vectorial & 32 B & costosa &
  Es homomórfico y ocultante, pero abrir una sola posición no es barato. \\
\addlinespace
KZG & 48 B & 48 B, constante &
  Setup de confianza y emparejamientos. A cambio, abre varias posiciones con una
  sola prueba. \\
\bottomrule
\end{tabularx}
\end{table}

\begin{nota}{Qué se paga en cada caso}
Merkle no pide confianza a nadie y paga con pruebas logarítmicas. KZG da pruebas
constantes y pide una ceremonia. No hay una respuesta correcta: hay dos
compromisos de ingeniería distintos.
\end{nota}

% ===========================================================================
\section{Dónde se usa esto de verdad}
% ===========================================================================

Un compromiso rara vez es el producto final: casi siempre es la primera mitad de
un protocolo en dos tiempos. El patrón \emph{comprometer ahora, revelar después}
aparece en sitios que parecen no tener nada que ver entre sí.

\subsection{Pruebas de conocimiento cero}

Todo protocolo sigma tiene la misma coreografía de tres pasos, y el primero es un
compromiso. Es lo que impide que el probador adapte su respuesta al desafío.

\begin{figure}[htbp]
\centering
\begin{tikzpicture}[
  font=\small,
  act/.style={draw=inkc!55,fill=inkc!8,rounded corners=2pt,
              minimum width=22mm,minimum height=6mm},
  >={Stealth[length=2mm]}
]
\node[act] (P) at (0,0) {Probador};
\node[act] (V) at (9.5,0) {Verificador};
\draw[inkc!35,dashed] (P) -- ++(0,-3.9);
\draw[inkc!35,dashed] (V) -- ++(0,-3.9);

\draw[->,violetdeep,thick] (0,-1.0) -- node[above,font=\scriptsize,violetdeep]
  {1. compromiso $a$} (9.5,-1.0);
\draw[<-,amberc,thick] (0,-2.1) -- node[above,font=\scriptsize,amberc]
  {2. desafío $e$ al azar} (9.5,-2.1);
\draw[->,mintc,thick] (0,-3.2) -- node[above,font=\scriptsize,mintc]
  {3. respuesta $z$} (9.5,-3.2);

\node[anchor=north west,align=left,font=\scriptsize,inkc,text width=118mm]
  at (-1.2,-4.3)
  {El compromiso del paso 1 fija la aleatoriedad del probador \emph{antes} de que
   conozca $e$. Sin él podría elegir $a$ después de ver el desafío y responder
   correctamente sin saber el secreto: la prueba dejaría de probar nada.};
\end{tikzpicture}
\caption{Protocolo sigma. El compromiso es lo que convierte tres mensajes en una
prueba.}
\label{fig:sigma}
\end{figure}

\subsection{Subastas a sobre cerrado}

El problema de pujar en público es que el último en hablar gana con la mínima
ventaja, sin importar cuánto valoraba el bien. Con compromisos, todos publican
$c = H(\text{puja} \concat r)$ en la fase 1 y abren en la fase 2. Quien intente
revelar otra cifra queda descalificado automáticamente: el hash no cuadra.

\begin{table}[htbp]
\centering\small
\caption{La misma subasta, con y sin compromisos. Cuatro postores; Dave es el
último en hablar.}
\label{tab:subasta}
\begin{tabular}{@{}l c c c@{}}
\toprule
\textbf{Postor} & \textbf{Valoración real} & \textbf{Sin compromiso} &
\textbf{Con compromiso} \\
\midrule
Alice & 120 & puja 120 & puja 120 \\
Bob   & 185 & puja 185 & puja 185 \\
Carol & 150 & puja 150 & puja 150 \\
Dave  & 170 & \textcolor{coralc}{puja 186} & intenta revelar 125 \\
\midrule
\textbf{Resultado} & --- &
  \textcolor{coralc}{\textbf{gana Dave con 186}} &
  \textcolor{mintc}{\textbf{gana Bob con 185}} \\
 & & {\scriptsize vio las pujas ajenas} & {\scriptsize Dave, descalificado} \\
\bottomrule
\end{tabular}
\end{table}

\begin{laboratorio}{Laboratorio 2.7 — Subasta a sobre cerrado}
Ejecuta la subasta de la tabla~\ref{tab:subasta} en las tres variantes: honesta,
con Dave haciendo trampa, y sin compromisos.

\textbf{Qué observar.} El compromiso no impide mentir: hace que \textbf{mentir
sea detectable}. Dave puede revelar la cifra que quiera; lo que no puede es que
el hash cuadre.
\end{laboratorio}

\subsection{Aleatoriedad pública}

Sorteos, asignación de turnos de validador, loterías en cadena: si cada
participante publica el compromiso de un número y luego todos abren, nadie pudo
elegir su aportación viendo la de los demás. El resultado es la suma. Es un
volado con $n$ jugadores, y es la base de los esquemas commit-reveal de RANDAO.

\begin{aviso}{El anti-patrón que hay que reconocer}
Un commit-reveal donde alguien puede \textbf{negarse a abrir} sin coste no es
justo: el último en revelar decide si participa según le convenga, y con ello
sesga el resultado. La solución habitual es exigir un depósito que se pierde si
no se abre.
\end{aviso}

\begin{nota}{En producción}
\textbf{Bitcoin}: la raíz de Merkle del encabezado es un compromiso a la lista de
transacciones. \textbf{Ethereum}: RANDAO usa commit-reveal para la aleatoriedad
del consenso y EIP-4844 usa KZG para los blobs. \textbf{Liquid y Grin}: Pedersen
para los montos, Bulletproofs para las pruebas de rango.
\textbf{Certificate Transparency}: un árbol de Merkle público como compromiso a
todos los certificados emitidos. \textbf{ENS}: registro en dos fases para que
nadie se adelante al ver tu solicitud.
\end{nota}

% ===========================================================================
\section{Síntesis: cuál elegir, y qué se rompe si falla}
% ===========================================================================

\begin{table}[htbp]
\centering\small
\caption{Los cuatro esquemas del apunte. La columna que decide casi siempre es
la última: \emph{qué necesitas hacer con el sobre cerrado}.}
\label{tab:sintesis}
\begin{tabularx}{\textwidth}{@{}l l l l >{\raggedright\arraybackslash}X@{}}
\toprule
\textbf{Esquema} & \textbf{Ocultamiento} & \textbf{Vinculación} &
\textbf{Tamaño} & \textbf{Qué permite} \\
\midrule
$H(m \concat r)$ \newline {\scriptsize SHA-256} &
  computacional & computacional & 32 B &
  Nada más que abrir. Sin setup, sin supuestos algebraicos, resistente a
  cuántico. \\
\addlinespace
Pedersen \newline {\scriptsize $g^{m}h^{r}$} &
  \textcolor{mintc}{\textbf{perfecto}} & dlog & 32 B &
  Sumar y restar sobres cerrados. Base de las transacciones confidenciales. \\
\addlinespace
Merkle \newline {\scriptsize vectorial} &
  \textcolor{coralc}{ninguno$^{*}$} & colisión & 32 B \newline $+\,32\log n$ &
  Abrir una posición sin revelar el resto. $^{*}$Oculta sólo si las hojas llevan
  sal. \\
\addlinespace
KZG \newline {\scriptsize polinómico} &
  computacional & q-SDH $+$ setup & 48 B \newline $+\,48$ B &
  Abrir cualquier evaluación con prueba constante. Exige ceremonia de confianza. \\
\bottomrule
\end{tabularx}
\end{table}

\begin{figure}[htbp]
\centering
\begin{tikzpicture}
\begin{axis}[
  xbar, width=11.2cm, height=5.4cm,
  bar width=11pt,
  xlabel={alcance del daño},
  xmin=0, xmax=115,
  xtick={0,25,50,75,100},
  xticklabels={,la ronda,,el protocolo,},
  symbolic y coords={{Prueba de rango},{Setup honesto},{Vinculación},{Ocultamiento}},
  ytick=data,
  yticklabel style={font=\footnotesize,align=right},
  grid=major, grid style={inkc!12},
  tick label style={font=\footnotesize}, label style={font=\small},
  enlarge y limits=0.22,
  nodes near coords, nodes near coords align={horizontal},
  point meta=explicit symbolic,
  every node near coord/.append style={font=\scriptsize,color=inkc}
]
\addplot[fill=skyc!45,draw=skyc!80!black] coordinates {
  (72,{Ocultamiento}) [subastas, votos, sorteos]};
\addplot[fill=coralc!55,draw=coralc!80!black] coordinates {
  (100,{Vinculación}) [todo el protocolo]};
\addplot[fill=amberc!55,draw=amberc!80!black] coordinates {
  (100,{Setup honesto}) [la vinculación, en silencio]};
\addplot[fill=violetc!45,draw=violetc!80!black] coordinates {
  (86,{Prueba de rango}) [inflación indetectable]};
\end{axis}
\end{tikzpicture}
\caption{Qué se rompe cuando falla cada pieza. Perder el \emph{ocultamiento}
arruina la ronda en curso, y se nota. Perder la \emph{vinculación} arruina el
protocolo entero y, si viene de un setup envenenado, no deja ni rastro.}
\label{fig:dano}
\end{figure}

\subsection{Una regla de decisión en tres preguntas}

\begin{enumerate}[leftmargin=2.2em,itemsep=0.4em]
  \item \textbf{¿Necesitas operar con el valor cerrado?} Si la respuesta es sí
    ---sumar montos, agregar votos---, es Pedersen. Si es no, el hash es más
    barato y no pide nada a cambio.
  \item \textbf{¿El secreto debe sobrevivir décadas?} Entonces quieres
    ocultamiento perfecto, porque un adversario futuro con más cómputo no cambia
    nada. Pedersen otra vez.
  \item \textbf{¿Te comprometes a muchas cosas?} Merkle si no quieres depender de
    ninguna ceremonia; KZG si necesitas pruebas de tamaño constante y puedes
    justificar el setup.
\end{enumerate}

\begin{nota}{Nota cuántica}
El logaritmo discreto cae ante el algoritmo de Shor. Eso significa que la
\textbf{vinculación} de Pedersen y de KZG no sobrevive a un computador cuántico
relevante, mientras que su ocultamiento perfecto sí: no se apoya en ningún
supuesto. El compromiso de hash está en la situación inversa y más cómoda: Grover
sólo divide a la mitad los bits de seguridad, así que SHA-256 conserva del orden
de 128.

La lectura práctica: para un \emph{registro} que deba seguir siendo inalterable
dentro de treinta años, un compromiso de hash es hoy la apuesta conservadora.
Para un \emph{secreto} que deba seguir siendo secreto dentro de treinta años, es
Pedersen.
\end{nota}

\begin{aviso}{Lo que hay que llevarse}
Un compromiso convierte una promesa en un objeto verificable. Las dos propiedades
apuntan a adversarios opuestos, no pueden ser ambas perfectas, y la que se rompe
en silencio ---la vinculación, vía un setup envenenado--- es siempre la más
peligrosa.
\end{aviso}

\clearpage
% ===========================================================================
\appendix
\section{Resumen de fórmulas}
% ===========================================================================

\begin{table}[htbp]
\centering\small
\caption{Todo el módulo en nueve renglones.}
\begin{tabular}{@{}l l l@{}}
\toprule
\textbf{Concepto} & \textbf{Expresión} & \textbf{Dónde} \\
\midrule
Compromiso de hash        & $c = H(m \concat r)$ & §2 \\
Cota del ocultamiento     & $\Adv^{\text{ocult}} \le q|\mathcal{M}|/2^{\lambda}$ & §2 \\
Cota de la vinculación    & $\Adv^{\text{vinc}} \le \Adv^{\text{col}}_H$ & Prop.~\ref{prop:hashbind} \\
Compromiso de Pedersen    & $C(m,r) = g^{m}h^{r}$ & §4 \\
Ruptura $\Rightarrow$ dlog & $\log_g h = (m-m')(r'-r)^{-1}$ & Prop.~\ref{prop:pedbind} \\
Homomorfismo              & $C(a,r_1)C(b,r_2) = C(a{+}b, r_1{+}r_2)$ & §5 \\
Balance confidencial      & $\prod C_{\text{ent}} = \prod C_{\text{sal}} \cdot C_{\text{com}}$ & §5 \\
Equivocación con trampa   & $r' = r + (m-m')x^{-1}$ & Lema~\ref{lema:trapdoor} \\
Prueba KZG                & $\pi = g^{(f(\tau)-y)/(\tau-z)}$ & §7 \\
\bottomrule
\end{tabular}
\end{table}

% ===========================================================================
\section{Ejercicios propuestos}
% ===========================================================================

\begin{enumerate}[leftmargin=2.4em,itemsep=0.5em]
  \item Un estudiante propone $c = H(r \concat m)$ en vez de $H(m \concat r)$.
    ¿Cambia algo respecto a la vinculación? ¿Y si $H$ es de Damgård--Merkle y el
    atacante controla la longitud de $m$? (Pista: módulo 01, extensión de
    longitud.)

  \item Demuestra que el compromiso $c = g^{m}$ ---Pedersen sin el término
    $h^{r}$--- es \emph{perfectamente vinculante} y no oculta nada cuando
    $\mathcal{M}$ es pequeño. Sitúalo en el cuadrante de la
    figura~\ref{fig:cuadrante}.

  \item En la demostración de la Proposición~\ref{prop:pedbind} se afirma que
    $r \neq r'$. Explica por qué el argumento se rompe si $q$ no es primo, y
    construye un contraejemplo con $q = 12$.

  \item Una transacción confidencial tiene una entrada de valor $v$ y dos salidas
    de valores $v_1, v_2$ con $v_1 + v_2 = v$. Escribe la condición que deben
    cumplir los factores de cegado para que el nodo acepte, y explica por qué el
    emisor puede elegirlos libremente salvo uno.

  \item Usando el Lema~\ref{lema:trapdoor}, calcula a mano el $r'$ que abre en
    $m' = 7$ un compromiso hecho con $m = 3$, $r = 5$, en el grupo $q = 1013$,
    $p = 2027$, $g = 4$, sabiendo que $h = g^{11}$. Verifica el resultado.

  \item Un protocolo commit-reveal de aleatoriedad tiene $n$ participantes y el
    resultado es la suma de los valores abiertos. Si el último puede negarse a
    abrir sin penalización, ¿cuántos resultados distintos puede forzar? ¿Y si
    pueden coaligarse $k$ de ellos?

  \item Compara el coste total en bytes de probar la pertenencia de $t$
    elementos distintos de un conjunto de $2^{20}$, con Merkle y con KZG. ¿A
    partir de qué $t$ conviene cada uno? Usa la figura~\ref{fig:tamanos}.
\end{enumerate}

\vfill
\begin{center}
\color{inkc}\small
\rule{0.35\textwidth}{0.4pt}\\[0.6em]
\textsc{ChainNotes} · Módulo 02 de 9 · Anterior: funciones hash · Siguiente: firma digital\\
{\footnotesize Versión LaTeX del apunte interactivo \texttt{temas/compromisos.html}.
Todas las figuras se generan con TikZ y pgfplots.}
\end{center}

\end{document}
